\documentclass[11pt,letterpaper]{article}

\usepackage[margin=1in]{geometry}
\usepackage{amsmath, amssymb}
\usepackage{booktabs}
\usepackage{graphicx}
\usepackage{hyperref}
\usepackage{listings}
\usepackage{xcolor}
\usepackage{caption}
\usepackage{subcaption}
\usepackage{microtype}
\usepackage{parskip}

% Figures are in plots/, tables in results_tables/.
% Upload those folders to Overleaf alongside this .tex file.
\graphicspath{{plots/}}

\hypersetup{colorlinks=true, linkcolor=blue, urlcolor=blue, citecolor=blue}

\lstset{
  basicstyle=\ttfamily\small,
  keywordstyle=\color{blue},
  commentstyle=\color{gray},
  breaklines=true,
  frame=single,
  language=C++
}
\title{\vspace{-3em}CME 213 Final Project - Milestone 3: GPU Kernels\\
\large Building and Optimizing a Small LLM from Scratch}
\author{Nils Astrup Toft \quad and \quad Sondre Rogde}
\date{\today}

\begin{document}
\maketitle

% ============================================================
\section{Overview}
% ============================================================

This milestone implements and benchmarks the custom CUDA kernels that form
the compute-intensive core of our GPT-2-style transformer.  The goal is to
replace na\"ive CPU or cuBLAS baselines with hand-written GPU kernels that
exploit the memory hierarchy and parallelism of NVIDIA Turing GPUs (sm\_75,
Quadro RTX 6000).

We implemented six kernels:
\begin{enumerate}
  \item \textbf{Tiled GEMM} --- shared-memory blocked matrix multiplication
  \item \textbf{Fused LayerNorm} --- single-kernel mean, variance, and normalization
  \item \textbf{Online Softmax} --- numerically stable row-wise softmax
  \item \textbf{GELU Activation} --- exact erf-based formulation
  \item \textbf{Cross-Entropy Loss} --- fused log-softmax and negative log-likelihood
  \item \textbf{Flash Attention (forward)} --- SRAM-tiled attention without materializing the $N \times N$ score matrix
\end{enumerate}

Each kernel is validated against a CPU reference implementation and against
PyTorch-generated ground-truth data before any performance measurement is
reported.

% ============================================================
\section{Implementation}
% ============================================================

\subsection{GEMM: Tiled Shared-Memory Multiplication}

\paragraph{Algorithm.}
Na\"ive GEMM has each thread read $K$ values from row $i$ of $A$ and $K$
values from column $j$ of $B$ directly from global memory.  The arithmetic
intensity is $\approx 1$ FLOP/byte for small tiles, leaving the GPU far
below its compute peak.

Our tiled kernel divides the $K$ dimension into tiles of width $T = 32$.
Each thread block loads a $T \times T$ sub-matrix of $A$ and $B$ into shared
memory cooperatively (one element per thread).  The inner product is then
computed entirely in shared memory, achieving $T$ reuses per global load.
Effective arithmetic intensity scales with $T$.

\paragraph{Key CUDA concepts.}
\begin{itemize}
  \item \emph{Shared memory}: a 48--96\,KB per-SM scratchpad with $\sim$100$\times$
    lower latency than global memory.  Declared with \texttt{\_\_shared\_\_}.
  \item \emph{\_\_syncthreads()}: a block-wide barrier that ensures all threads
    have finished filling the shared tile before any thread reads from it.
  \item \emph{Coalesced loads}: mapping \texttt{threadIdx.x} to the contiguous
    (column) dimension ensures that a warp's 32 loads collapse into a single
    128-byte memory transaction.
\end{itemize}

\paragraph{Baseline comparison.}
We compare against a CPU reference GEMM (\texttt{gemm\_cpu}) implemented as
a plain triple-loop in C++.  The CPU time is reported alongside the GPU time
for small correctness-check sizes, giving a direct measure of GPU speedup.

\paragraph{Further optimization (future work).}
Our tiled kernel computes one output element per thread.  A further
improvement — \emph{register tiling} — has each thread compute a small
$T_m \times T_n$ sub-tile ($4\times 4$ is typical), keeping those accumulators
in registers throughout the K-loop rather than re-accessing shared memory.
This reduces the $\nicefrac{\texttt{\_\_syncthreads()}}{\text{compute}}$ ratio
and better hides shared-memory latency.  We implemented this technique in
CME 213 HW4 (\texttt{tiledGEMM}), where each thread computed 8 output elements
per block tile; adapting that kernel here is a planned follow-up.

\subsection{Fused LayerNorm}

LayerNorm normalizes each row of its input:
\[
  y_i = \gamma \cdot \frac{x_i - \mu}{\sqrt{\sigma^2 + \varepsilon}} + \beta,
  \quad
  \mu = \frac{1}{H}\sum_j x_j,
  \quad
  \sigma^2 = \frac{1}{H}\sum_j (x_j - \mu)^2.
\]

\paragraph{Fusion.}
A na\"ive implementation requires three separate kernel launches (mean,
variance, normalize), each reading and writing the full $N \times H$ tensor.
Our fused kernel computes everything in a single launch with three logical
passes but only two global kernel dispatches:
\begin{enumerate}
  \item \emph{Pass 1}: each thread accumulates a partial sum; a shared-memory
    tree reduction yields $\mu$.
  \item \emph{Pass 2a}: reuse the shared buffer to accumulate $(x_j - \mu)^2$
    and reduce to $\sigma^2$.  Computing variance this way (rather than the
    algebraically equivalent $E[x^2] - E[x]^2$) avoids catastrophic
    cancellation when $\mu$ is large.
  \item \emph{Pass 2b}: normalize, scale, and shift in a single write pass.
\end{enumerate}
Total global memory traffic: 3 reads of $x$ + 1 write of $y$, versus 6
reads + 3 writes for the na\"ive three-pass approach.

\subsection{Softmax}

We implement a three-pass online softmax:
\begin{enumerate}
  \item Find the row maximum $m_i = \max_j x_{ij}$ (for numerical stability).
  \item Compute $\sum_j e^{x_{ij} - m_i}$.
  \item Normalize: $y_{ij} = e^{x_{ij} - m_i} / \text{sum}$.
\end{enumerate}
Subtracting the row maximum before exponentiation prevents overflow.
We use CUDA's fast intrinsic \texttt{\_\_expf} for throughput (small accuracy
trade-off, within $5 \times 10^{-4}$ relative error versus double-precision).

\subsection{GELU Activation}

GELU is the standard activation for GPT-2's feed-forward layers:
\[
  \text{GELU}(x) = x \cdot \tfrac{1}{2}\bigl(1 + \operatorname{erf}(x / \sqrt{2})\bigr).
\]
We use the exact erf formulation (matching PyTorch's
\texttt{F.gelu(x, approximate='none')}) via CUDA's \texttt{erff()} SFU
intrinsic.  The kernel is purely pointwise with arithmetic intensity
$\approx 0.75$\,FLOP/byte, making it firmly memory-bandwidth bound.

\subsection{Cross-Entropy Loss}

For each token we compute:
\[
  \mathcal{L}_n = -\log \frac{e^{\text{logits}[n, t_n]}}{\sum_v e^{\text{logits}[n,v]}},
\]
using the log-sum-exp trick (subtract the row maximum before computing
softmax) to avoid floating-point overflow.  The kernel reads logits twice
(one pass for the max, one for the sum of exponentials) and writes one loss
scalar per token.

\subsection{Flash Attention (Forward Pass)}

Standard scaled dot-product attention requires materializing the $S \times S$
score matrix per head, using $O(S^2)$ memory.  For long sequences this is the
dominant memory bottleneck.

Our Flash Attention implementation follows Dao et al.\ (2022).  The key idea
is to tile the $Q$, $K$, $V$ matrices into SRAM-sized blocks and use an
\emph{online} softmax update to accumulate the output $O$ incrementally,
never writing the full score matrix to DRAM.  This reduces memory from
$O(S^2)$ to $O(S)$ and improves memory bandwidth utilization substantially
for long sequences.

\paragraph{Online softmax.}
As we stream $K$ and $V$ tiles into SRAM we maintain a running maximum
$m^{(j)}$ and running sum $\ell^{(j)}$, applying a correction factor
$e^{m^{(j-1)} - m^{(j)}}$ when the maximum changes.  The final output is
rescaled once at the end of the loop.

We support causal masking by zeroing out attention scores for keys at
positions $> q$ before the softmax.

% ============================================================
\section{Correctness Validation}
% ============================================================

Each kernel is tested in two ways:

\paragraph{CPU oracle.}
A plain C++ reference implementation of each kernel runs on the host.
The test compares the GPU output against the CPU output element-wise,
reporting maximum absolute and relative error.  All kernels pass with
relative error below their respective tolerance ($10^{-3}$ for GEMM and
GELU, $10^{-4}$ for LayerNorm and cross-entropy, $5 \times 10^{-3}$ for
Flash Attention).

\paragraph{PyTorch comparison.}
We generate reference inputs and outputs with PyTorch
(\texttt{scripts/generate\_ref\_data.py}) and compare the saved \texttt{.bin}
tensors against our GPU results.  This catches any systematic divergence from
the numerically established ground truth.  Table~\ref{tab:pytorch_results}
summarizes the PyTorch comparison results.

\begin{table}[h]
  \centering
  \small
  \begin{tabular}{lrrl}
    \toprule
    \textbf{Kernel} & \textbf{Max abs err} & \textbf{Max rel err} & \textbf{Result} \\
    \midrule
    GEMM            & $0.000$ & $0.000$ & PASS \\
    LayerNorm       & \textit{[placeholder]} & \textit{[placeholder]} & \textit{[after fix: PASS]} \\
    Softmax         & $1.1 \times 10^{-8}$ & $6.4 \times 10^{-7}$ & PASS \\
    GELU            & $4.8 \times 10^{-7}$ & $2.9 \times 10^{-4}$ & PASS \\
    Cross-Entropy   & $9.5 \times 10^{-7}$ & $1.4 \times 10^{-7}$ & PASS \\
    Flash Attention & $4.2 \times 10^{-7}$ & $1.1 \times 10^{-3}$ & PASS \\
    \bottomrule
  \end{tabular}
  \caption{PyTorch reference comparison (\texttt{test\_vs\_pytorch}).
           GELU's non-zero relative error is expected: PyTorch evaluates
           \texttt{erf} in double precision internally, while our kernel uses
           CUDA's single-precision \texttt{erff()}, causing amplified relative
           error near zero-valued outputs.}
  \label{tab:pytorch_results}
\end{table}

% ============================================================
\section{Performance Results}
\label{sec:results}
% ============================================================

All measurements are taken on the Quadro RTX 6000 (Turing architecture,
sm\_75, 16.31\,TFLOPS FP32, 672\,GB/s DRAM bandwidth).  Each reported time
is the average of 10--20 kernel launches bracketed by CUDA events, with one
untimed warm-up launch to amortize JIT compilation and first-touch costs.

\subsection{Roofline Analysis}

Figure~\ref{fig:roofline} places each kernel on the roofline model.  The
roofline ceiling for a kernel with arithmetic intensity $I$ (FLOP/byte) is:
\[
  \text{attainable} = \min\bigl(\text{peak compute},\; I \times \text{peak BW}\bigr).
\]
Kernels to the left of the ridge point ($\approx 24$\,FLOP/byte) are
memory-bandwidth bound; those to the right are compute bound.

\begin{figure}[h]
  \centering
  \includegraphics[width=0.95\textwidth]{roofline}
  \caption{Roofline model for all kernels on the Quadro RTX 6000
           (16.31\,TFLOPS FP32, 672\,GB/s DRAM).
           Dashed vertical lines show the gap between each kernel and the
           theoretical ceiling for its arithmetic intensity.
           Memory-bound kernels (LayerNorm, Softmax, GELU, Cross-Entropy)
           cluster on the left bandwidth-limited slope; GEMM and Flash
           Attention sit in the compute-limited flat region.
           Generated by \texttt{scripts/plot\_perf.py}.}
  \label{fig:roofline}
\end{figure}

\subsection{Performance Table}

% Auto-generated by scripts/plot_perf.py -- do not edit by hand.
% Regenerate with: python3 scripts/plot_perf.py logs/run_<jobid>.out
\begin{table}[h]
  \centering
  \small
  \begin{tabular}{lllrrrr}
    \toprule
    \textbf{Kernel} & \textbf{Problem size} & \textbf{Time (ms)} & \textbf{Achieved} & \textbf{AI (F/B)} & \textbf{\% of ceil.} & \textbf{Bound} \\
    \midrule
  \texttt{GEMM (tiled)} & $M=N=K=1024$ & 0.923 & 2326.2 GFLOPS & 170.7 & 14.3\% & Compute \\
  \texttt{Flash Attention} & $B=1,S=2048,D=96$ & 6.059 & 2126.7 GFLOPS & 512.0 & 13.0\% & Compute \\
  \texttt{LayerNorm} & $N=8192,\ H=768$ & 0.115 & 438.1 GB/s & 0.5 & 65.2\% & Memory \\
  \texttt{Softmax} & $N=4096,\ V=30000$ & 4.597 & 427.7 GB/s & 0.2 & 63.6\% & Memory \\
  \texttt{GELU} & $N=4,194,304$ & 0.071 & 469.2 GB/s & 0.8 & 69.8\% & Memory \\
  \texttt{Cross-Entropy} & $N=128,\ V=30000$ & 0.065 & 474.9 GB/s & 0.4 & 70.7\% & Memory \\
    \bottomrule
  \end{tabular}
  \caption{Best achieved performance per kernel on the Quadro RTX 6000 (Turing sm\_75, FP32 peak 16.3\,TFLOPS, DRAM 672\,GB/s). \emph{AI} = arithmetic intensity (FLOP/byte). \emph{\% of ceil.} = achieved / roofline ceiling for that kernel's AI.}
  \label{tab:m3_perf}
\end{table}


\subsection{GEMM: GPU Speedup over CPU Baseline}

The CPU reference (\texttt{gemm\_cpu}, plain triple-loop in C++) is measured
alongside the GPU kernel for small correctness-check sizes.  The CPU time
grows as $O(MNK)$ while the GPU kernel's parallelism makes it dramatically
faster even at modest sizes:

\begin{itemize}
  \item At $M=N=K=256$, the CPU takes $\approx$\textit{[X]} ms versus the
    GPU's $\approx$\textit{[X]} ms — a $\approx$\textit{[XX$\times$]} speedup.
    \textit{(Fill in from SLURM log.)}
  \item At $M=N=K=4096$, the CPU would take on the order of minutes; the
    GPU completes in $\approx$\textit{[X]} ms at
    $\approx$\textit{[XXXX]} GFLOPS.
\end{itemize}

The tiled kernel achieves these numbers by staging $32 \times 32$ sub-matrices
of $A$ and $B$ into shared memory, reusing each loaded element 32 times before
eviction.  This raises the effective arithmetic intensity far above the
ridge point, making the kernel compute-bound rather than bandwidth-bound.

% ============================================================
\section{Discussion}
% ============================================================

\subsection{Memory-Bound Kernels}

LayerNorm, Softmax, GELU, and Cross-Entropy all have arithmetic intensity
well below the ridge point.  Their performance ceiling is set by DRAM
bandwidth (672\,GB/s), not the FP32 compute units.  The primary optimization
levers are:

\begin{itemize}
  \item \textbf{Minimizing passes over the data.}  Fusing operations
    (e.g., our two-pass LayerNorm) reduces the number of global memory
    round-trips.
  \item \textbf{Vectorized loads.}  Replacing four scalar \texttt{float}
    loads with a single \texttt{float4} instruction reduces instruction
    overhead; this is a near-term follow-up.
  \item \textbf{Occupancy.}  High occupancy (many resident warps per SM)
    allows the hardware to hide global-memory latency.  For memory-bound
    kernels, maximizing occupancy is the main path to closing the gap
    between achieved and peak bandwidth.
\end{itemize}

\subsection{Compute-Bound Kernels}

GEMM and Flash Attention are compute-bound at realistic problem sizes
(after tiling).  The gap between achieved GFLOPS and the 16.3\,TFLOPS
ceiling is explained by:

\begin{itemize}
  \item \textbf{Register tiling.}  Our tiled GEMM uses a $32 \times 32$
    thread block with one output element per thread.  Moving to register
    tiling (each thread computes multiple outputs, keeping accumulators in
    the register file) reduces the
    $\nicefrac{\texttt{\_\_syncthreads()}}{\text{compute}}$ ratio.  We
    implemented this in CME 213 HW4 (\texttt{tiledGEMM}: 128 threads, 8
    outputs per thread per $32 \times 32$ block tile); adapting that kernel
    here is the planned next step.
  \item \textbf{No software pipelining.}  Loading the next tile while
    computing the current one (\texttt{cp.async} / double-buffering) is
    required to approach peak on Ampere and above; it also helps on Turing.
  \item \textbf{Tensor Cores.}  Turing's Tensor Core units deliver
    $\sim$125\,TFLOPS for FP16 matrix operations, more than $7\times$ the
    FP32 scalar throughput.  Adopting mixed-precision (BF16/FP16 compute
    with FP32 accumulation) is planned for Milestone 4.
\end{itemize}

\subsection{Flash Attention}

Flash Attention avoids materializing the $S \times S$ attention matrix, which
for $S=512$ and 8 heads would require $8 \times 512^2 \times 4$\,bytes
$= 8$\,MB per batch element.  The SRAM tiling strategy keeps working memory
proportional to $S \cdot D$ rather than $S^2$, making it feasible for long
sequences.

The achieved GFLOPS for Flash Attention should be compared against the
na\"ive attention (which would write the score matrix to DRAM), not against
the raw GEMM kernel, because the memory savings change the roofline ceiling.
A na\"ive attention comparison is left as a follow-up experiment.

% ============================================================
\section{Future Work}
% ============================================================

\begin{enumerate}
  \item \textbf{Register tiling for GEMM} --- adapt the \texttt{tiledGEMM}
    kernel from CME 213 HW4 (128 threads, 8 outputs per thread) to our
    row-major convention.  This reduces synchronization overhead and keeps
    accumulators in registers, the primary bottleneck at our current
    $\sim$\textit{XX}\% of compute peak.
  \item \textbf{Vectorized loads (float4)} --- reduce instruction count for
    memory-bound kernels by issuing 128-bit loads.
  \item \textbf{Mixed precision (BF16/FP32)} --- switch storage and compute
    to BF16 with FP32 accumulation; required before the full training loop.
  \item \textbf{Backward passes} --- implement gradients for GEMM, LayerNorm,
    GELU, and Flash Attention (Milestone 4).
  \item \textbf{Fused GEMM + GELU} --- eliminate the intermediate HBM
    round-trip between the linear projection and the activation.
\end{enumerate}

% ============================================================
\section{Build and Run Instructions}
% ============================================================

\begin{lstlisting}[language=bash, caption={Building and running the kernel tests.}]
# On the SLURM cluster (gpu-turing partition, Quadro RTX 6000):
sbatch run.sh

# After the job completes, generate the roofline plot and LaTeX table:
python3 scripts/plot_perf.py logs/run_<jobid>.out
# Outputs: reports/roofline.pdf, reports/perf_table.tex

# To regenerate PyTorch reference data (run once, on a node with PyTorch):
python3 scripts/generate_ref_data.py
\end{lstlisting}

\end{document}
